\newpage
\section{Clases de Stiefel-Whitney}

\cite[\S 4]{ClasesCaracteristicas}.

Damos inicio a esta sección introduciendo los axiomas que caracterizan a las clases de Stiefel-Whitney. Es posible establecer tanto existencia como unicidad de las clases de cohomología que satisfacen estos axiomas, y aunque es algo que complementaría de buena forma a este informe, no es algo que trabajaremos durante el mismo. Para tal exposición, referirse a las secciones 7 y 8 de \cite{ClasesCaracteristicas}. Durante esta sección, asumimos estos hechos.

\begin{axiom}
	A cada fibrado vectorial $\xi$ le corresponde una secuencia de clases de cohomología
	\[ w_{i}(\xi) \in H^{i}(B(\xi); \Z \slash 2\Z), \quad i = 0,1,2,\dots \]
	llamadas las \textbf{clases de Stiefel-Whitney de $\xi$}. La clase $w_{0}(\xi)$ es igual a la unidad
	\[ 1 \in H^{0}(B(\xi); \Z \slash 2\Z), \]
	y $w_{i}(\xi)$ es igual a cero para $i > n$ si $\xi$ es un fibrado de $n$-planos.
\end{axiom}

\begin{axiom}[Naturalidad]
	Si $f \colon B(\xi) \to B(\eta)$ está cubierto por un mapa de fibrados de $\xi$ a $\eta$, entonces
	\[ w_{i}(\xi) = f^{*}w_{i}(\eta). \]
\end{axiom}

\begin{axiom}[Teorema del Producto de Whitney]
	Si $\xi$ y $\eta$ son fibrados sobre un mismo espacio base, entonces 
	\[ w_{k}(\xi \oplus \eta) = \sum_{i=0}^{k} w_{i}(\xi) \smile w_{k-i}(\eta). \]
\end{axiom}

\begin{axiom}
	Para el fibrado de linea $\gamma_{1}^{1}$ sobre $\P^{1}$, la clase de Stiefel-Whitney $w_{1}(\gamma_{1}^{1})$ es no nula.
\end{axiom}

Notemos que el axioma 2 puede ser reescrito de la siguiente manera
\[ w_{i}(f^{*}\eta) = f^{*}w_{i}(\eta), \]
donde, por un lado, $f^{*}$ representa el fibrado inducido por $f$ y, por otro, el mapa inducido en cohomología. En cierto sentido, estamos estableciendo que las clases conmutan con los pullback's en los contextos apropiados. \par

Por otro lado, y a modo de ejemplificar el comportamiento dado por el axioma 3, notar que
\begin{align*}
	w_{1}(\xi \oplus \eta) &= w_{1}(\xi) + w_{1}(\eta) \\
	w_{2}(\xi \oplus \eta) &= w_{2}(\xi) + w_{1}(\xi)w_{1}(\eta) + w_{2}(\eta) \\
	w_{3}(\xi \oplus \eta) &= w_{3}(\xi) + w_{2}(\xi)w_{1}(\eta) + w_{1}(\xi)w_{2}(\eta) + w_{3}(\eta).
\end{align*}
Tal como en el ejemplo, normalmente omitiremos el uso del símbolo $\smile$ para el producto cup.

\subsection{Consecuencias de los Axiomas}

\begin{lemma}
	Si $\xi$ es isomorfo a $\eta$, entonces $w_{i}(\xi) = w_{i}(\eta)$.
\end{lemma}
\begin{proof}
	...
\end{proof}

\begin{lemma}
	Si $\varepsilon$ es un fibrado vectorial trivial, entonces $w_{i}(\varepsilon) = 0$ para $i > 0$.
\end{lemma}
\begin{proof}
	...
\end{proof}

\begin{lemma}
	Si $\varepsilon$ es trivial, entonces $w_{i}(\varepsilon \oplus \eta) = w_{i}(\eta)$.
\end{lemma}
\begin{proof}
	...
\end{proof}

\begin{lemma}
	Si $\xi$ es un fibrado de $n$-planos con una métrica Euclidiana que posee una sección cruzada no nula, entonces $w_{n}(\xi) = 0$. Aún más, si $\xi$ posee $k$ secciones cruzadas linealmente independientes, entonces
	\[ w_{n-k+1}(\xi) = w_{n-k+2}(\xi) = \cdots = w_{n}(\xi) = 0. \]
\end{lemma}
\begin{proof}
	...
\end{proof}

Un caso particular del axioma 3 que es interesante revisar es cuando la suma de Whitney $\xi \oplus \eta$ es trivial. Por el lemma 5.2, se cumplen las siguientes relaciones
\begin{align*}
	& w_{1}(\xi) + w_{1}(\eta) = 0 \\
	& w_{2}(\xi) + w_{1}(\xi) w_{1}(\eta) + w_{2}(\eta) = 0 \\
	& w_{3}(\xi) + w_{2}(\xi) w_{1}(\eta) + w_{1}(\xi) w_{2}(\eta) + w_{3}(\eta) = 0 \\
	& \vdots
\end{align*}
De esta manera, obtenemos una manera inductiva de expresar $w_{i}(\eta)$ en términos de las clases de Stiefel-Whitney de $\xi$. Es decir, si conocemos las clases de $\xi$, es fácil calcular las clases de $\eta$ en este caso. \par

Ahora, introducimos un formalismo que será clave para argumentar ciertas cosas en lo que procede.

\begin{definition}
	$H^{\Pi}(B; \Z \slash 2\Z)$ denotará el anillo que consiste de todas las series formales infinitas
	\[ a = a_{0} + a_{1} + a_{2} + \cdots \]
	donde $a_{i} \in H^{i}(B; \Z \slash 2\Z)$. El producto en este anillo está dado por la fórmula
	\begin{align*}
		& (a_{0} + a_{1} + a_{2} + \cdots)(b_{0} + b_{1} + b_{2} + \cdots) \\
		& \quad = (a_{0}b_{0}) + (a_{1}b_{0} + a_{0}b_{1}) + (a_{2}b_{0} + a_{1}b_{1} + a_{0}b_{2}) + \cdots. \\
	\end{align*}
	Este producto es conmutativo, puesto que estamos trabajando mod $2$, y es asociativo. Aditivamente, $H^{\Pi}(B; \Z \slash 2\Z)$ se comporta como el producto cartesiano de los grupos $H^{i}(B; \Z \slash 2\Z)$.
\end{definition}

\begin{definition}
	La \textbf{clase total de Stiefel-Whitney} de un fibrado de $n$ planos $\xi$ sobre $B$ se define por el elemento 
	\[ w(\xi) = 1 + w_{1}(\xi) + w_{2}(\xi) + \cdots + w_{n}(\xi) + 0 + \cdots \]
	en $H^{\Pi}(B; \Z \slash 2\Z)$.
\end{definition}

Notemos entonces que, el producto de Whitney (axioma 3), ahora puede ser expresado por la fórmula
\[ w(\xi \oplus \eta) = w(\xi) w(\eta). \]
Es claro que esta igualdad nos está diciendo que los elementos coinciden en cada grado del anillo de cohomología, dado que estamos trabajando con sumas formales.

\begin{lemma}
	La colección de todas las series infinitas
	\[ w = 1 + w_{1} + w_{2} + \cdots \in H^{\Pi}(B; \Z \slash 2\Z) \]
	forma un grupo conmutativo con la multiplicación.
\end{lemma}
\begin{proof}
	...
\end{proof}

Ahora, consideremos dos fibrados vectoriales $\xi$ y $\eta$ sobre el mismo espacio base. Ya establecimos que podemos escribir
\[ w(\xi \oplus \eta) = w(\xi)w(\eta), \]
y gracias al lemma anterior, podemos reescribir esto de la siguiente forma
\[ w(\eta) = \overline{w}(\xi) w(\xi \oplus \eta). \]
En particular, si la suma de Whitney es trivial, entonces
\[ w(\eta) = \overline{w}(\xi). \]
Esto implica un caso especial de suma importancia.

\begin{theorem}[Dualidad de Whitney]\label{Dualidad de Whitney}
	Si $\tau_{M}$ es el fibrado tangente de una variedad $M$ dentro de un espacio euclídeo y $\nu$ es el fibrado normal asociado, entonces
	\[ w_{i}(\nu) = \overline{w}_{i}(\tau_{M}). \]
\end{theorem}

Con estas herramientas, buscamos calcular las clases de Stiefel-Whitney en algunos casos especiales. Además, con frecuencia utilizaremos la abreviación $w(M)$ para la clase total de Stiefel-Whitney del fibrado tangente $\tau_{M}$.

\begin{example}
	La clase total de Stiefel-Whitney del fibrado tangente de la esfera unitaria $S^{n}$ es igual a $1$. Esto nos dice que el fibrado tangente de $S^{n}$ y el fibrado trivial sobre $S^{n}$ no puede ser distinguido mediante clases de Stiefel-Whitney.
\end{example}
\begin{proof}
	...
\end{proof}

Durante el resto de esta sección nos centraremos principalmente en los fibrados asociados a los espacios proyectivos $\P^{n}$. Recordemos, de la sección anterior, que el grupo $H^{i}(\P^{n}; \Z \slash 2\Z)$ es el grupo cíclico de orden $2$ para $0 \leq i \leq n$ y es cero para valores mayores de $i$. Aún más, si $a$ denota el elemento no nulo de $H^{1}(\P^{n}; \Z \slash 2\Z)$, entonces el anillo de cohomología de $\P^{n}$ está dado por
\[ H^{*}(\P^{n}; \Z \slash 2\Z) \cong (\Z \slash 2\Z) \slash (a^{n+1}). \]

% Mencionar que el mapa f^{*} \colon H^{n}(\P^{n}, \Z \slash 2\Z) \to H^{n}(S^{n}, \Z \slash 2\Z) cumple que f^{*}(a^{n}) = (f^{*}a)^{n}, aunque esto es algo que quiero mencionar en la sección de cohomología, cuando mencione el lemma que relaciona el producto cup con los mapas inducidos.

\begin{example}
	La clase total de Stiefel-Whitney del fibrado tautológico $\gamma_{n}^{1}$ sobre $\P^{n}$ está dado por
	\[ w(\gamma_{n}^{1}) = 1 + a. \]
\end{example}
\begin{proof}
	...
\end{proof}

\begin{example}
	Por definición, el fibrado tautológico $\gamma_{n}^{1}$ sobre $\P^{n}$ está contenido como sub-fibrado en el fibrado trivial $\varepsilon^{n+1}$. Luego, sea $\gamma^{\perp}$ el complemento ortogonal de $\gamma_{n}^{1}$ en $\varepsilon^{n+1}$. Entonces,
	\[ w(\gamma^{\perp}) = 1 + a + a^{2} + \cdots + a^{n}. \]
\end{example}
\begin{proof}
	...
\end{proof}

Utilizando algunas de las propiedades que hemos enunciado hasta este momento, somos capaces de calcular $w(\P^{n})$.

\begin{theorem}\label{Clase del Tangente}
	La suma de Whitney $\tau \oplus \varepsilon^{1}$ es isomorfo a la suma de Whitney de $n+1$ fibrados canónicos, $\gamma_{n}^{1} \oplus \cdots \oplus \gamma_{n}^{1}$. Luego, la clase total de Stiefel-Whitney de $\P^{n}$ está dada por
	\[ w(\P^{n}) = (1 + a)^{n+1} = 1 + \binom{n+1}{1} a + \binom{n+1}{2} a^{2} + \cdots \binom{n+1}{n} a^{n}. \]
\end{theorem}
\begin{proof}
	...
\end{proof}

El siguiente corolario nos da una obstrucción a la posibilidad de paralelizar el fibrado tangente de cada espacio proyectivo.

\begin{corollary}
	La clase $w(\P^{n})$ es igual a $1$ si y solo si $n+1$ es una potencia de $2$. Por lo tanto, los único espacios proyectivos que pueden ser paralelizables son $\P^{1},\P^{3},\P^{7},\P^{15},\dots$.
\end{corollary}
\begin{proof}
	...
\end{proof}

He aquí una ocasión en que las clases de Stiefel-Whitney son capaces de darnos información, pero terminan no siendo completamente determinantes. En efecto, se puede ver que $\P^{1}, \P^{3}$ y $\P^{7}$ son paralelizables. En cambio, los espacios proyectivos de mayor dimensión $\P^{15},\P^{31},\dots$ no son paralelizables. \\[0.1pt]
% esto es posible citarlo apropiadamente %

Similarmente, uno podría preguntarse si $\gamma^{\perp}$ isomorfo al fibrado tangente $\tau$ de $\P^{n}$, para todo $n$. Sabemos que cuando $n = 2$ esto es cierto, pero las clases de Stiefel-Whitney terminan siendo una buena obstrucción a la pregunta. \par En efecto, notemos que en caso que sean isomorfos, dado que $\gamma_{n}^{1} \oplus \gamma^{\perp} \cong \varepsilon^{n+1}$, entonces
\[ w(\gamma_{n}^{1} \oplus \tau) = 1. \]
Ahora bien, también se cumple que
\[ w(\gamma_{n}^{1} \oplus \tau) = w(\gamma_{n}^{1}) w(\tau) = (1+a)^{n+2}. \]
Pero $(1+a)^{n+2} = 1$ si y sólo si $n+2$ es una potencia de $2$, por el mismo argumento dado en la demostración del corolario. Por lo tanto, las únicas dimensiones donde es posible que el isomorfismo en cuestión ocurra son $\P^{2},\P^{6},\P^{14},\P^{30},\dots$. \par


\subsection{Inmersiones}

Con tal de seguir ejemplificando la naturaleza de obstrucción de las clases de Stiefel-Whitney, damos una aplicación final del teorema \ref{Clase del Tangente}. Preguntémonos entonces qué espacios proyectivos pueden ser inmersos en un espacio Euclídeo de una dimensión particular. \par
Para ello, notemos primero que si una variedad $M$ de dimensión $n$ puede ser inmersa en $\R^{n+k}$, el teorema de dualidad de Whitney (\ref{Dualidad de Whitney}) nos dice que
\[ w_{i}(\nu) = \overline{w}_{i}(M), \]
y esto implica que las clases de Stiefel-Whitney duales $\overline{w}_{i}(M)$ son cero para $i > k$. Es decir, para que el espacio pueda potencialmente ser inmerso en $\R^{n+k}$, se debe cumplir esta condición sobre los duales.

\begin{example}
	En el caso del espacio proyectivo $\P^{9}$, dado que
	\[ w(\P^{9}) = (1 + a)^{10} = 1 + a^{2} + a^{4}, \]
	entonces
	\[ \overline{w}(\P^{9}) = 1 + a^{2} + a^{4} + a^{6}. \]
	Es decir, si $\P^{9}$ puede ser inmerso en $\R^{n + k}$, entonces $k$ debe ser por lo menos $6$.
\end{example}

El caso más impresionante en $\P^{n}$ es cuando $n$ es una potencia de $2$, y es lo que vemos en el siguiente teorema.

\begin{theorem}
	Si $\P^{2^{r}}$ puede ser inmerso en $\R^{2^{r} + k}$, entonces $k \geq 2^{r} - 1$.
\end{theorem}
\begin{proof}
	...
\end{proof}
